\section{Related Literature and Contribution}\label{sec:literature}

The financing question posed in Section~\ref{sec:introduction} draws on five strands of research, each addressing a different link between a technological gain and a recurrent social commitment. The relevant distinction is between what each strand disciplines and what it leaves unresolved. Read in sequence, the literatures identify the uncertainty in the aggregate shock, its allocation across factor incomes, its conversion into public revenue, and the cost of the policy to be financed.

\paragraph{Macroeconomic effects of AI and automation.}
Macroeconomic assessments span task accounting, micro-to-macro aggregation, and evidence on technology diffusion. Syntheses that organize AI's productivity, distributional, and growth channels do not collapse them into a single aggregate path \citep{oecd2024impactai}. Task-based models begin with the activities that technology can perform and aggregate from their economic weight, the associated cost savings, and the creation or reinstatement of tasks \citep{acemoglu2018race,acemoglu2019automation,acemoglu2024simple}. This approach makes occupational exposure an input rather than an aggregate growth estimate. A complementary micro-to-macro literature combines task-level performance, exposure, adoption, and propagation across sectors to construct aggregate productivity paths, including frontier applications under high exposure and adoption \citep{filippucci2024miracle,oecd2025g7}. Field and experimental studies document productivity gains in particular work settings and professional writing tasks \citep{brynjolfsson2023generative,noy2023experimental}. Cross-country evidence on earlier technologies documents substantial heterogeneity in adoption timing across technologies and economies \citep{comin2010diffusion}. Complementary intangible investment can also delay the appearance of gains in measured productivity, producing a productivity J-curve rather than immediate realization \citep{brynjolfsson2021jcurve}. Preparedness-based analysis therefore treats digital infrastructure, skills, innovation linkages, and institutions as conditions for converting exposure into use \citep{cazzaniga2024genai}. These studies discipline task-level gains, diffusion, and timing, but do not by themselves identify aggregate output or the tax bases produced by diffusion. The aggregation exercises, in turn, depend on assumptions about adoption and transmission beyond the settings in which performance is observed \citep{acemoglu2024simple,filippucci2024miracle}. This strand therefore narrows plausible mechanisms without selecting a single magnitude, diffusion speed, or duration for the aggregate shock. Those open margins motivate a scenario input rather than a point forecast.

\paragraph{Distributional incidence and factor allocation.}
Automation can change both output and its allocation. In the task framework, displacement transfers activities away from labor, while reinstatement creates new labor-intensive tasks; productivity and labor income therefore need not move proportionately \citep{acemoglu2018race,acemoglu2019automation}. Regional labor-market mapping distinguishes exposure associated with productivity-enhancing transformation from automation risk and identifies the digital divide as a constraint on realizing the former \citep{worldbankilo2024lacgenai}. The wider factor-share literature likewise treats the division between labor and nonlabor income as an empirical object rather than a fixed accounting constant \citep{karabarbounis2014global}. Measurement adds a second distinction. Labor income must be adjusted for self-employment and mixed income, so the national-accounts residual outside measured employee compensation is not synonymous with corporate profit \citep{gollin2002getting}. It may contain mixed income, depreciation, operating surplus, and rents with different legal and fiscal treatments. An aggregate AI gain consequently does not reveal how much enters wages, taxable business income, consumption, or a separately reachable rent base. This literature establishes why factor allocation matters, but leaves the allocation of a particular projected increment across taxable bases unresolved.

\paragraph{Fiscal capacity and effective capture.}
Research on AI and public finance describes fiscal policy as a channel through which technological gains can be distributed and captured, using social protection, labor-market policy, and labor, capital, or corporate taxation \citep{brollo2024fiscal,korinek2026public}. Optimal-tax analyses of automation further show that the treatment of capital that substitutes for labor depends on the transition being considered and on the available tax instruments \citep{guerreiro2022robots}. Whether a proposed instrument produces revenue also depends on state capacity. \citet{besley2014taxation} relate tax collection in developing countries to investments in fiscal institutions, while \citet{jensen2022employment} links the expansion of employee taxation to the information generated by formal employment and third-party reporting. Work on tax capacity similarly distinguishes the reach of a tax system from the size of the underlying economy \citep{gaspar2016tax}. Panel estimates of tax buoyancy show that short- and long-run revenue responses differ by tax category and country group, so aggregate income growth need not map one-for-one into receipts \citep{dudine2018buoyancy}. That evidence disciplines historical capture comparisons, but does not identify the marginal response to an AI increment. These mechanisms create a wedge between economic income, a statutory rate, and realized marginal revenue. Behavioral base erosion, enforcement, and cross-border profit shifting belong in that mapping when corporate income is part of the base; reviews of international corporate tax avoidance document both multiple shifting channels and measurement blind spots \citep{beer2020international,brollo2024fiscal}. Nor is additional revenue alone a complete measure of fiscal space. Standard debt dynamics link the debt path to primary balances and growth-interest conditions, so a revenue-cost crossing does not establish debt consistency by itself \citep{escolano2010debt,imf2016fiscalspace}. The open object is thus effective, debt-consistent capture from a specified income increment, not the tax rate in isolation.

\paragraph{Social-protection design in economies with large informal sectors.}
Universal basic income, guaranteed minimum income, social pensions, and targeted transfers impose different eligibility, information, and budget requirements. Reviews of UBI therefore treat it as a policy architecture with explicit financing and interaction with existing benefits, rather than as a generic response to technological change \citep{gentilini2020ubi,oecd2017basicincome}. In developing-country settings, the comparison between universal and targeted transfers turns partly on imperfect information: proxy-means tests reduce coverage but introduce targeting errors when income is difficult to observe \citep{hanna2018universal}. Informal and non-employee work also limits the information available to tax and transfer systems \citep{jensen2022employment}. Informality is also associated with small firm scale, credit constraints, and weak complementary investment, linking the transfer-information problem to productive adoption and fiscal reach \citep{worldbank2021informality}. From the financing side, universal social-protection floors require recurrent fiscal space rather than a one-time resource envelope \citep{cattaneo2024financing}. Policy-modeling work on labor-saving technologies compares basic income with alternatives such as job guarantees and working-time reduction, showing that program architecture shapes distributional, employment, and fiscal effects \citep{dalessandro2025exploring}. This strand clarifies what governments may wish to finance and why targeting changes cost. It does not determine whether an uncertain AI-related revenue flow can cover the recurring cost of any one instrument.

\paragraph{AI--UBI feasibility and the policy-modeling tradition.}
The closest threshold studies already ask financing questions. \citet{nayebi2025threshold} derives an AI-capability threshold for a rent-funded UBI in the United States, and \citet{ito2026ubi} simulates UBI feasibility in Japan with dynamic social-security offsets. These are feasibility exercises rather than static incidence studies. Their common unit, however, is one country and one instrument, so neither is designed to compare the same financing condition across alternative social-protection commitments and fiscal architectures. Within the \emph{Journal of Policy Modeling} literature, \citet{caamal2022ubi} use tax-benefit microsimulation to evaluate UBI schemes in Mexico, while \citet{lara2026basic} evaluate a basic-income--flat-tax reform in Spain through poverty, inequality, and efficiency outcomes. \citet{dalessandro2025exploring} broaden that policy dialogue by comparing several responses to labor-saving technologies in a dynamic modeling environment. These studies evaluate specified reforms and their incidence, efficiency, or budget implications. They do not jointly ask how the same conditional financing source performs across countries, instruments, technological scenarios, and fiscal regimes.

Taken together, the reviewed strands provide the components of the financing problem but not a common interface among them. This paper supplies a scenario-based fiscal-feasibility test that keeps the external technological gain conditional, separates domestic translation from factor allocation and effective capture, prices alternative instruments on a common denominator, and screens an accounting crossing for debt consistency. It also inverts the test to identify the required gain, translation, or capture and the persistence horizon associated with the financing condition. Modularity is substantive here: a change in technological evidence need not be recast as a tax reform, and a change in fiscal architecture need not alter the technology scenario. The comparison complements incidence, welfare, and program-design analysis; it does not replace them or turn the scenario inputs into causal estimates. Section~\ref{sec:framework} now converts these open margins into separate analytical objects and derives the common threshold.
